% ===== SECTION 7: COMPANY PRACTICES DISCLOSURE =====
\section{\texorpdfstring{\color{MidnightBlue}Company Practices Disclosure}{Company Practices Disclosure}}

\textit{``Please disclose your company practices in the following areas:''}

\begin{sidebar}{RFI Reference}
Section C: Questions to Field --- Company Practices Disclosure (Items 1--3)
\end{sidebar}

% ----- Item 1 -----
\subsection{Item 1: Driver Recruitment Pipeline}

Sentry approaches driver recruitment from two complementary angles: as a technology platform that equips bus operators with workforce management tools, and as an organization with direct hiring experience managing vehicle operators across our portion of the MTA Access-A-Ride paratransit network in New York City. In that program, Sentry coordinates a provider network of transportation companies and has developed proven systems for tracking applicant pipelines, credential status, and time-to-road-ready metrics at scale. We bring that operational discipline to RIDE's Statewide Transportation System, providing zone operators with the digital infrastructure to recruit, credential, and retain qualified drivers in one of the most competitive labor markets in the Northeast.

\begin{description}[style=nextline, leftmargin=0.5cm]
    \item[\faIcon{bullhorn} \textbf{Recruitment Technology:}] Sentry's platform includes an integrated applicant tracking system (ATS) that manages the full recruitment funnel---from initial application through background check, CDL training enrollment, credentialing, and road-ready certification. The system supports targeted digital advertising campaigns on job boards and social media, manages referral bonus programs with automated payout tracking, and integrates directly with CDL training school enrollment systems so that applicant progress is visible in real time. Operators can see at a glance how many candidates are at each pipeline stage and where bottlenecks are forming, enabling proactive intervention before driver shortages affect service.

    \item[\faIcon{graduation-cap} \textbf{CDL Pipeline Development:}] Sentry's workforce management module supports ``earn while you learn'' CDL pipeline programs by tracking candidates from initial interest through classroom instruction, behind-the-wheel training, CDL examination, and school bus endorsement. The platform manages partnerships with CDL training schools, scheduling candidates into available cohorts and tracking completion rates. For Rhode Island specifically, the system is configured to track the RI ``white card'' credentialing process---including the multi-day course at CCRI and DMV driving tests---as well as the new temporary license provision for CDL holders residing in Massachusetts and Connecticut. This ensures that border-state drivers are routed through the abbreviated training pathway while their full RI credentials are processed, expanding the available driver pool without compromising compliance.

    \item[\faIcon{handshake} \textbf{Retention as Recruitment:}] The most cost-effective recruitment strategy is retention. Sentry's platform includes scheduling fairness algorithms that distribute route assignments equitably, route preference systems that allow experienced drivers to bid on preferred routes, and performance recognition dashboards that track on-time rates, safety records, and attendance for merit-based incentives. Real-time communication tools---including in-app messaging, push notifications, and automated shift reminders---keep drivers informed and engaged. In our MTA Access-A-Ride operation, where we coordinate scheduling for over 10,000 daily trips, these retention tools have contributed to a driver turnover rate significantly below the industry average. When drivers feel their schedules are fair, their performance is recognized, and communication is transparent, the constant churn of recruitment diminishes---a critical advantage in Rhode Island, where drivers are actively recruited away by Uber, DoorDash, RIPTA, FedEx, and Amazon.
\end{description}

\mybox{
\textbf{\faIcon{info-circle} Rhode Island Context:} The September 2024 service crisis demonstrated the consequences of inadequate driver recruitment---Dattco deployed CT-based drivers who lacked the RI ``white card'' (a multi-day course at CCRI plus DMV driving tests unique to Rhode Island). Even with the new temporary license provision for border-state CDL holders, recruitment pipelines must account for this RI-specific credentialing timeline. Additionally, driver competition from Uber, DoorDash, RIPTA, FedEx, and Amazon has intensified the shortage. RIDE's October~2024 system overview identifies this as the system's most significant threat. Sentry supports RIDE's proposed \textbf{Bus Driver Apprenticeship program} concept and can provide the workforce management technology to administer it at scale.
}{MidnightBlue}

% ----- Item 2 -----
\subsection{Item 2: Training Hours Per Driver}

Sentry's training management platform administers a structured, two-tier training program for all drivers and aides in the statewide system. Initial training provides comprehensive preparation before a driver operates a route independently, covering CDL preparation, defensive driving, student management, ADA compliance, technology systems, and emergency procedures. Annual refresher training reinforces critical skills, incorporates lessons learned from incident reviews, and ensures compliance with RIDE's 17-hour annual training requirement and the 8-hour de-escalation requirement for special education personnel. All training completion, certification status, and credential expiration dates are tracked digitally, with automated alerts sent to supervisors 30, 14, and 7 days before any credential expires.

\vspace{0.3cm}

\begin{tabularx}{\textwidth}{|X|c|c|}
\hline
\rowcolor{tableHeader}
\textcolor{white}{\textbf{Training Module}} &
\textcolor{white}{\textbf{Initial Hours}} &
\textcolor{white}{\textbf{Annual Refresher}} \\
\hline
CDL / School Bus Endorsement Preparation & 10 hrs & N/A \\ \hline
Defensive Driving (Smith System) & 8 hrs & 2 hrs \\ \hline
Special Education Student Handling & 6 hrs & 2 hrs \\ \hline
Wheelchair Securement \& ADA Compliance & 4 hrs & 2 hrs \\ \hline
Student Management \& De-Escalation & 8 hrs & 4 hrs \\ \hline
Technology Systems (Driver App, GPS, Cameras) & 4 hrs & 2 hrs \\ \hline
Emergency Evacuation Procedures & 4 hrs & 2 hrs \\ \hline
FERPA \& Student Privacy & 2 hrs & 1 hr \\ \hline
\rowcolor{tableHeader}
\textcolor{white}{\textbf{Total}} &
\textcolor{white}{\textbf{46 hrs}} &
\textcolor{white}{\textbf{15 hrs}} \\
\hline
\end{tabularx}

\vspace{0.3cm}

All training hours meet or exceed Rhode Island regulatory requirements, including the 10-hour school bus driver training course mandated by state law (or the abbreviated course now available for applicants who already hold a valid CDL, per the 2025 temporary license legislation). The 8-hour de-escalation module for special education drivers and aides exceeds RIDE's required hours, reflecting the heightened behavioral management demands on routes serving students with disabilities. Sentry's training management module delivers these courses through a blended format---online modules for knowledge-based content and in-person sessions for hands-on skills---and maintains a complete digital record of every training hour completed by every driver and aide in the system.

\vspace{0.3cm}

\mybox{
\textbf{\faIcon{info-circle} RIDE's Current Training Standards:} RIDE currently requires \textbf{17~hours of annual training} for all statewide drivers and aides, plus \textbf{8~hours of de-escalation training} specifically for drivers and aides assigned to special education buses. Sentry's training management module tracks certification status, schedules refresher courses, and alerts supervisors when drivers approach credential expiration---ensuring 100\% compliance with these requirements across all zone operators.
}{MidnightBlue}

% ----- Item 3 -----
\subsection{Item 3: Maintenance Intervals}

Sentry's fleet management platform enforces a tiered preventive maintenance (PM) schedule aligned with industry best practices and manufacturer specifications. The system integrates with GPS tracking hardware to track both mileage and calendar intervals for every vehicle in the fleet, automatically generating work orders when a vehicle approaches a PM threshold. Critically, the platform prevents dispatch of any vehicle that has exceeded its PM due date or has an open safety defect---ensuring that maintenance is never deferred due to operator scheduling pressure or driver shortages. This technology-enforced compliance is particularly important in a multi-operator environment where RIDE needs assurance that all zone contractors maintain their fleets to uniform standards.

\vspace{0.3cm}

\begin{tabularx}{\textwidth}{|X|c|X|}
\hline
\rowcolor{tableHeader}
\textcolor{white}{\textbf{Service Type}} &
\textcolor{white}{\textbf{Interval}} &
\textcolor{white}{\textbf{Scope}} \\
\hline
PM-A (Basic) & Every 6,000 miles or 60 days & Oil and filter change, all fluid levels, belt and hose inspection, tire pressure and tread depth, all interior and exterior lights, wiper blades, mirrors \\ \hline
PM-B (Intermediate) & Every 12,000 miles or 6 months & PM-A scope plus brake measurement and adjustment, suspension and steering component inspection, exhaust system check, battery load test, air system inspection, wheelchair lift functional test \\ \hline
PM-C (Comprehensive) & Every 36,000 miles or annually & PM-B scope plus transmission service, coolant system flush, full electrical system diagnostic, body and structural integrity inspection, seat frame and restraint inspection, complete HVAC service \\ \hline
Annual DOT Inspection & Annual & Full Federal Motor Carrier Safety Administration (FMCSA) annual inspection per 49~CFR~396.17, including all Rhode Island state-specific school bus inspection items; results recorded digitally and available to RIDE on demand \\ \hline
Pre-Trip Inspection & Daily & Driver walk-around per FMCSA \S 396.13---brakes, tires, lights, mirrors, emergency exits, wheelchair lift, fluid leaks, and student safety equipment---digitally logged via the Sentry driver app with GPS-stamped completion time \\ \hline
Post-Trip Inspection & Daily & Mandatory interior sweep for remaining students (child-check system with rear-of-bus button), equipment condition assessment, and digital defect reporting with photo upload capability for immediate maintenance dispatch \\ \hline
\end{tabularx}

\vspace{0.3cm}

\begin{sidebar}{Technology-Enforced Compliance}
Sentry's fleet management module is designed to integrate with GPS tracking hardware to monitor real-time mileage accumulation and calendar intervals for every vehicle. The system automatically generates work orders at 90\% of each PM threshold, escalates alerts to zone supervisors and RIDE's dashboard at 95\%, and \textbf{blocks dispatch of any vehicle that exceeds its PM due date or carries an unresolved safety defect}. This ``hard stop'' mechanism ensures that no operator can defer maintenance to keep a vehicle in service during driver shortages or peak demand periods. RIDE receives a weekly fleet compliance scorecard showing PM completion rates, open defects, and average time-to-repair for every zone operator---providing the independent oversight necessary to maintain uniform safety standards across the entire statewide system.
\end{sidebar}
